\documentclass[12pt,a4paper]{article}
\usepackage[utf8]{inputenc}
\usepackage[T1]{fontenc}
\usepackage[english]{babel}
\usepackage{lmodern}
\usepackage[a4paper,margin=2.5cm]{geometry}
\usepackage{setspace}
\usepackage{hyperref}
\usepackage{enumitem}
\usepackage{graphicx}

\setstretch{1.1}
\hypersetup{
    colorlinks=true,
    linkcolor=black,
    urlcolor=blue
}

\title{Binary Classification of PCG Signals with a Focus on Preprocessing}
\author{
    \textbf{Artificial Intelligence from Engineering to Arts - a.y. 2025/2026}
    \\
    Gianluca Farinaccio - 548013 \\
    Alexis Osvaldo Hidalgo Martinez - 549484 \\
    \\
    \url{https://github.com/osvaldomar144/pcg-heart-sound-pipelines.git}
}

\begin{document}
\maketitle

\begin{abstract}
This project addresses binary classification of phonocardiograms (PCG) into
healthy/unhealthy classes, using data derived from the PhysioNet/CinC 2016 challenge.
The main objective is to systematically evaluate the impact of preprocessing before
training a CNN. Modular pipelines based on trim, band-pass filtering, and RMS
normalization were compared, followed by Mel transform and logarithmic compression.
The results show that preprocessing has a substantial impact on the trade-off between
overall accuracy and the ability to detect the minority unhealthy class. The main
contribution of this work is therefore methodological: a reproducible comparison between
pipelines and an analysis of the limitations of a baseline CNN on an imbalanced dataset.
\end{abstract}

\section{Introduction}
Automatic classification of heart sounds is a relevant problem for early screening and
triage, especially in low-resource settings. However, PCG signals are difficult to model
because they are often affected by environmental noise, acquisition artifacts, and
intra/inter-patient physiological variability. In addition, class imbalance can shift model
behavior toward the majority class.

In this project, an end-to-end pipeline was built, including configurable audio
preprocessing, conversion into a time-frequency representation (Mel spectrogram), and CNN
classification. Priority was given to preprocessing because, in practice, PCG models are
highly sensitive to how the signal is filtered, normalized, and temporally standardized.

\section{State of the art}
The PhysioNet/CinC 2016 challenge remains the historical benchmark most commonly used for
normal/abnormal PCG classification. In general, the literature reports significantly
better results than those obtained in this baseline project, but these performances almost
always rely on more complex pipelines.

State-of-the-art solutions are not based on a simple CNN applied to Mel spectrograms: they
typically integrate advanced preprocessing, explicit heart-cycle segmentation, and more
sophisticated noise-reduction strategies. In addition, beyond standard CNNs, sequential and
hybrid models are often employed, such as LSTM/BiLSTM architectures or combinations of
multiple architectural blocks.

In summary, the comparison with the state of the art suggests that this work is a solid
methodological starting point, but bridging the performance gap requires more advanced
preprocessing and modeling.

\section{Methods and technologies}
\subsection{Dataset}
The dataset was initially available as a single training block. To build a complete
experimental protocol, this set was split into train, validation, and test partitions using
a stratified 80\%-15\%-5\% split, keeping the healthy/unhealthy distribution as consistent as
possible across the three subsets.

The raw data consists of \texttt{.wav} audio files sampled at 2000 Hz, with variable duration
across recordings. The dataset is also heterogeneous with respect to acquisition context:
auscultation sites may differ from file to file, but this information is not provided in a
structured way in the metadata used in this project. The \texttt{unhealthy} class is also
aggregated: it indicates the presence of an anomaly, but does not systematically specify
the underlying cardiac pathology.

From a statistical perspective, the problem remains imbalanced. The Imbalance Ratio, defined
as \(N(\text{healthy}) / N(\text{unhealthy})\), is approximately 3.87: in practice, the
healthy class is almost four times more represented than unhealthy. This means that a naive
classifier always predicting healthy can still achieve apparently high accuracy, while being
clinically uninformative.

During training, this imbalance tends to favor the majority class during optimization, with
the risk of lower recall on the minority class and more false negatives on unhealthy.
Therefore, accuracy alone is not sufficient: evaluation should also include macro-F1,
class-wise recall, and confusion-matrix analysis.

\subsection{Implemented preprocessing pipelines}
Pipelines are defined in \texttt{configs/pipelines.json} with shared global parameters
(\texttt{target\_sr=2000}, \texttt{chunk\_seconds=8.13}, \texttt{n\_fft=512},
\texttt{hop\_length=128}, \texttt{n\_mels=128}, Mel band 20--300 Hz). The three variants
used in the comparison are \texttt{trim}, \texttt{bandpass}, and
\texttt{bandpass\_norm\_rms}.

In all pipelines, audio is loaded in mono and then temporally standardized with
\texttt{trim} to the target duration. In the \texttt{bandpass} variant, a Butterworth
band-pass filter (20--300 Hz, order 4, with \texttt{filtfilt}) is applied to attenuate
out-of-band components. In \texttt{bandpass\_norm\_rms}, RMS normalization is additionally
applied to reduce volume variability across recordings. The final representation is obtained
with \texttt{mel}, followed by \texttt{log1p} transformation to compress spectrogram dynamics
and stabilize the network input.

Pipeline parameters were not selected arbitrarily, but derived from iterative signal analysis
and constraints imposed by CNN input requirements. In particular, filter band, chunk
duration, and time-frequency parameters were chosen to balance general robustness,
preservation of useful information, and reduced risk of overfitting to this specific dataset.

With the selected parameters (\texttt{n\_mels=128}, \texttt{chunk\_seconds=8.13},
\texttt{hop\_length=128}), the output of each pipeline is a 128x128 log-Mel spectrogram,
stored in grayscale.

\paragraph{Qualitative healthy/unhealthy example}
To qualitatively visualize preprocessing effects, Figures \ref{fig:trim-log-pair} and
\ref{fig:bnr-log-pair} show the same healthy/unhealthy file pair after the \texttt{log} step,
for the \texttt{trim} and \texttt{bandpass\_norm\_rms} pipelines, respectively. In both
images, the left panel is healthy and the right panel is unhealthy.

\begin{figure}[htbp]
    \centering
    \includegraphics[width=\textwidth]{sanity/pipeline_trim_20260213_153137/04_log.png}
    \caption{Healthy/unhealthy pair with \texttt{trim} pipeline (log-Mel output)}
    \label{fig:trim-log-pair}
\end{figure}

\begin{figure}[htbp]
    \centering
    \includegraphics[width=\textwidth]{sanity/pipeline_bandpass_norm_rms_20260213_153230/06_log.png}
    \caption{Healthy/unhealthy pair with \texttt{bandpass\_norm\_rms} pipeline (log-Mel output)}
    \label{fig:bnr-log-pair}
\end{figure}

Visual analysis highlights a consistent difference between classes in both pipelines: the
healthy sample shows richer temporal structure and energy spread extending more into
mid-low bands, while the unhealthy sample shows energy more concentrated in low frequencies
and sparser temporal patterns. In the direct comparison between \texttt{trim} and
\texttt{bandpass\_norm\_rms}, global differences are moderate on this example, but the latter
produces slightly cleaner maps with more regular contrast, consistent with the combined
effect of band-pass filtering and RMS normalization.

\subsection{CNN architecture and training parameters}
The classifier is a ResNet18 initialized with pretrained ImageNet weights and adapted to
binary classification by replacing the final layer. Spectrograms are converted to grayscale
images replicated on 3 channels, resized to 224x224, and normalized.

Three training runs were performed, one for each pipeline described above. Training was run
for 10 epochs, batch size 32, Adam optimizer, and CrossEntropy loss. In addition,
\texttt{WeightedRandomSampler} is enabled because, as discussed above, the dataset is strongly
imbalanced, with an Imbalance Ratio of 3.87.

\subsection{Technology stack}
The project uses Python as scripting language,
\texttt{numpy}/\texttt{scipy}/\texttt{librosa}/\texttt{soundfile} for audio processing,
\texttt{matplotlib} for visualization and reporting, \texttt{torch} and \texttt{torchvision}
for training and evaluation, and \texttt{scikit-learn} for classification metrics. The whole
workflow is orchestrated through dedicated CLI scripts for preprocessing, training, sanity
checks, and test evaluation.

\section{Results}
Evaluation reports the three \texttt{model\_best} runs for pipelines \texttt{bandpass},
\texttt{bandpass\_norm\_rms}, and \texttt{trim}. All comparisons refer to the same test set
(324 samples: 258 healthy and 66 unhealthy).

\subsection*{Global metrics comparison}
Table~\ref{tab:metriche-globali} summarizes the main comparison across accuracy, precision,
recall, and weighted F1.

\begin{table}[htbp]
    \centering
    \caption{Global metrics comparison across pipelines}
    \label{tab:metriche-globali}
    \begin{tabular}{|l|c|c|c|c|}
        \hline
        \textbf{Pipeline} & \textbf{Accuracy} & \textbf{Precision weighted} & \textbf{Recall weighted} & \textbf{F1 weighted} \\
        \hline
        \texttt{bandpass} & 0.7531 & 0.7418 & 0.7531 & 0.7470 \\
        \hline
        \texttt{bandpass\_norm\_rms} & 0.7222 & 0.8461 & 0.7222 & 0.7489 \\
        \hline
        \texttt{trim} & 0.5926 & 0.7682 & 0.5926 & 0.6318 \\
        \hline
    \end{tabular}
\end{table}

Weighted F1 is especially useful in this context because it accounts for the actual class
frequencies in the test set, preventing the global result from being overly influenced by
low-support classes. At the same time, it remains more informative than plain accuracy
because it combines precision and recall into a single indicator.

The \texttt{bandpass} pipeline achieves the best accuracy, while
\texttt{bandpass\_norm\_rms} achieves the best weighted F1. The \texttt{trim} pipeline
remains the least effective globally.

\subsection*{Focus on the unhealthy class}
Since the unhealthy class is minority and clinically relevant,
Table~\ref{tab:metriche-unhealthy} reports class-specific precision, recall, and F1.

\begin{table}[htbp]
    \centering
    \caption{Metrics on the unhealthy class}
    \label{tab:metriche-unhealthy}
    \begin{tabular}{|l|c|c|c|}
        \hline
        \textbf{Pipeline} & \textbf{Precision unhealthy} & \textbf{Recall unhealthy} & \textbf{F1 unhealthy} \\
        \hline
        \texttt{bandpass} & 0.3793 & 0.3333 & 0.3548 \\
        \hline
        \texttt{bandpass\_norm\_rms} & 0.4143 & 0.8788 & 0.5631 \\
        \hline
        \texttt{trim} & 0.2963 & 0.7273 & 0.4211 \\
        \hline
    \end{tabular}
\end{table}

The most relevant value is unhealthy recall: \texttt{bandpass\_norm\_rms} sharply reduces
false negatives compared with the other pipelines, while \texttt{bandpass} favors the
healthy class and misses many unhealthy samples.

\subsection*{Visual support}
Figures~\ref{fig:cm-normalizzate} and \ref{fig:global-metrics-plots} show normalized
confusion matrices and global-metric plots for the three pipelines.

\begin{figure}[htbp]
    \centering
    \includegraphics[width=0.32\textwidth]{metrics/pipeline_banpass_model_best_test/confusion_matrix_normalized.png}
    \hfill
    \includegraphics[width=0.32\textwidth]{metrics/pipeline_bnr_model_best_test/confusion_matrix_normalized.png}
    \hfill
    \includegraphics[width=0.32\textwidth]{metrics/pipeline_trim_model_best_test/confusion_matrix_normalized.png}
    \caption{Normalized confusion matrices: \texttt{bandpass} (left),
    \texttt{bandpass\_norm\_rms} (center), \texttt{trim} (right)}
    \label{fig:cm-normalizzate}
\end{figure}

\begin{figure}[htbp]
    \centering
    \includegraphics[width=0.32\textwidth]{metrics/pipeline_banpass_model_best_test/global_metrics.png}
    \hfill
    \includegraphics[width=0.32\textwidth]{metrics/pipeline_bnr_model_best_test/global_metrics.png}
    \hfill
    \includegraphics[width=0.32\textwidth]{metrics/pipeline_trim_model_best_test/global_metrics.png}
    \caption{Global-metric plots: \texttt{bandpass} (left),
    \texttt{bandpass\_norm\_rms} (center), \texttt{trim} (right)}
    \label{fig:global-metrics-plots}
\end{figure}

In summary, results highlight two distinct behaviors: \texttt{bandpass} is the best choice
when maximizing overall accuracy, while \texttt{bandpass\_norm\_rms} is more robust when the
priority is detection of the unhealthy class (higher recall and better weighted F1).

\section{Conclusions}
The project produced a solid and reproducible experimental basis to compare preprocessing
pipelines for binary PCG classification. The main finding is that preprocessing strongly
changes the trade-off between specificity on healthy and sensitivity on unhealthy: there is
no single ``best'' pipeline in absolute terms, but a choice tied to the target metric and
application context.

A key methodological point is that, under class imbalance, accuracy and val\_acc can be
overly optimistic if not complemented by class-wise metrics, weighted F1, and error
analysis. Another important limitation is the amount of available data: the number of
samples, especially for the unhealthy class, reduces model generalization and increases
sensitivity to dataset-specific characteristics.

Future work should therefore prioritize multi-seed/cross-validation comparisons, targeted
data augmentation strategies to increase variability and improve class balance, and
validation of the whole pipeline on larger and more heterogeneous external datasets. In
parallel, it remains relevant to evaluate hybrid models with sequential/attention
components and calibrate decision thresholds based on the clinical cost of false negatives.

\section*{Web references (primary sources)}
\begin{enumerate}[leftmargin=*,label=\arabic*.]
    \item PhysioNet Challenge 2016 (official overview): \url{https://physionet.org/content/challenge-2016/}
    \item Zabihi et al., Heart Sound Classification from Unsegmented Phonocardiograms Using Deep Learning (CinC 2016): \url{https://www.cinc.org/archives/2016/pdf/236-175.pdf}
    \item Springer et al., Logistic Regression-HSMM-Based Heart Sound Segmentation (IEEE TBME 2016): \url{https://pubmed.ncbi.nlm.nih.gov/26340769/}
\end{enumerate}

\end{document}
